\section{Delivery summary}
\label{sec:conclusion}
The delivered MIPS program runs a two-player match from the twelve-piece opening to a win, while keeping board state, turn ownership and the displayed square numbers consistent. A useful acceptance case is the black capture sequence \code{11x2 2x9 9x18}: the first jump promotes the pawn, the second uses its new backward direction, and the player is not asked to choose another starting piece between jumps. It exercises input, board updates and the rules together.

The optional settings build on that base without silently changing the assignment rules. Red can rotate the display while history retains one numbering system. A player can request hints without requiring the first capture, or enable forced capture as a separate rule. The bitmap view reads the same position as the terminal rather than maintaining a second game.

The five saved classic-mode MARS transcripts match their expected outputs after line-ending and surrounding-whitespace normalization; Black wins games 001, 003 and 005, and Red wins 002 and 004. These games check complete sessions. The targeted positions and comparisons in the verification section address questions a complete game alone cannot settle, such as whether every legal continuation was generated.

The most visible implementation saving was avoiding unnecessary redraws: the recorded two-tile update went from 81,263 to 1,908 simulated instructions, while the mean across five classic games went from 101,776 to 84,812. Those are different workloads, not competing estimates of the same speedup. The deliverable is the playable game, its reproducible runs and the evidence above; the instruction counts explain the cost of specific choices on MARS, not hardware execution time.
